\subsection{The constructed symbolic polynomial normalizer}

\begin{hbdefinition}{def:p12-har-poly-regime}{Exact grammar, algebraic hypotheses and normal-form convention}
Use exactly the five-constructor arithmetic syntax T and E from
\cref{def:p12-har-autarkic-syntax,def:p12-har-autarkic-interpreter}:
Z, U, Var i, Add s t and Mul s t. The following program implements
the omitted source G using positive finite data and their generated
primitive recursion in the explicitly historical Set-based regime.
It is not a beta-only HOL term, an imported polynomial library or the
generic compiler cited elsewhere by the chapter.

For interpretation fix A:Set, typed zero/unit/addition/multiplication,
and primitive book equality on A. Assume the actual commutative-semiring
laws: associative and commutative addition with zero identity;
associative and commutative multiplication with unit identity;
both distributive laws; and zero annihilation on both sides.
These are explicit propositional equation families. They are supplied
by the constructed arithmetic library when A=nat. An arbitrary carrier
with uninterpreted operations does not have these laws.
Every well-typed valuation nu:nat-to-A in a valid context is admitted.
Neither DN, choice, functional extensionality nor parametricity is used.

Normal data are sparse polynomials: a monomial is a nondecreasing finite
word of variable indices, repetitions recording powers; coefficients
are positive naturals; rows have distinct words in strictly descending
lexicographic order, with the empty word smaller than every nonempty
word. The quotation convention is a right-associated sum with no
trailing zero, natural literals made from U/Add, coefficient-one factors
omitted, and no trailing unit in nonempty monomial products. This
explicit convention may differ from implicit source parenthesization;
the example below gives both the actual output and its association bridge
\cite[pp.~1221--1222]{BarendregtGeuvers2001}.
\end{hbdefinition}

\begin{hbdefinition}{def:p12-har-poly-data}{The actual finite data types and recursion support}
Declare the strictly positive data
\[
\begin{aligned}
 W&:\mathsf{Set},& []&:W,&i::w&:W\quad(i:\mathsf{nat},w:W),\\
 R&:\mathsf{Set},&\mathsf{row}&:\mathsf{nat}\to W\to R,\\
 P&:\mathsf{Set},&\epsilon&:P,&r\mathbin{\#}p&:P\quad(r:R,p:P).
\end{aligned}
\]
For a fixed admitted target C, \(\mathrm{Rec}_W\) has seed z:C and branch
nat-to-W-to-C-to-C; \(\mathrm{Rec}_R\) has a single branch nat-to-W-to-C;
\(\mathrm{Rec}_P\) has seed z:C and branch R-to-P-to-C-to-C. Their computations
at constructors are respectively z/f i w(rec w), f c w, and
z/f r p(rec p). Generated Prop induction receives exactly the
recursive-field hypotheses; R has none. These specified instances,
not an assumed list/map library or a top-sort-polymorphic constant,
are all finite-list support used below. A case operation is the same
recursor with its recursive value ignored.

For each fixed C put
\(I_C:=\lambda b:\mathsf{nat}.\lambda u,v:C.
\mathsf{nat\_rec}\ v(\lambda j:\mathsf{nat}.\lambda r:C.u)b\).
Thus \(I_C\) O u v contracts to v and \(I_C\) one u v to u.
Use the already constructed eqb, leb, ltb, andb and orb data codes.
Write \(\mathsf{empty}_W\)/\(\mathsf{empty}_P\) for their constructor cases returning one at
the empty constructor and O at a cons. Row coefficient/word and
word head/tail projections are their constructor cases; head []=O,
tail []=[] are explicit defaults.

All displayed functional equations below are abbreviations for these
fixed-target constructor recursions; branch types and nested recursion
orders are checked in the annex. No well-founded search is silently
postulated from an ordering relation.
\end{hbdefinition}

\begin{hbdefinition}{def:p12-har-poly-word-code}{Computed word equality, order and product}
The actual equality and strict-lexicographic codes are
\[
\begin{aligned}
 \mathsf{weq}([],[])&=\mathsf{one},&
 \mathsf{weq}([],j::v)&=O,&\mathsf{weq}(i::u,[])&=O,\\
 \mathsf{weq}(i::u,j::v)&=
          \mathsf{andb}(\mathsf{eqb}\ i\ j)(\mathsf{weq}(u,v)),\\
 \mathsf{wlt}([],[])&=O,&
 \mathsf{wlt}([],j::v)&=\mathsf{one},&\mathsf{wlt}(i::u,[])&=O,\\
 \mathsf{wlt}(i::u,j::v)&=
   \mathsf{orb}(\mathsf{ltb}\ i\ j)
       (\mathsf{andb}(\mathsf{eqb}\ i\ j)(\mathsf{wlt}(u,v))).
\end{aligned}
\]
Each is an outer \(\mathrm{Rec}_W\) at target W-to-nat, with a case on the
second word in its seed/step and recursive value q:W-to-nat used
only at v. This is an actual nested primitive recursion.
Let u<lex v be the inductive Prop relation with constructors:
[]<lex j::v; i<j implies i::u<lex j::v; and i=j together with
u<lex v implies i::u<lex j::v. The conclusion depends only on
data endpoints, not on a proof argument.

Monomial multiplication merges words:
\[
 \mathsf{MW}([],v)=v,\qquad \mathsf{MW}(i::u,[])=i::u,
\]
and on two conses it computes
\[
 \mathsf{MW}(i::u,j::v)=
 I_W(\mathsf{leb}\ i\ j)
       (i::\mathsf{MW}(u,j::v))
       (j::\mathsf{MW}(i::u,v)).
\]
To make both recursive calls legitimate, its outer \(\mathrm{Rec}_W\) has target
W-to-W, seed lambda v.v, and cons branch
\[
 \lambda i:\mathsf{nat}.\lambda u:W.\lambda q:W\to W.
 \lambda v:W.\mathrm{Rec}_W(i::u)
  (\lambda j:\mathsf{nat}.\lambda v':W.\lambda z:W.
     I_W(\mathsf{leb}\ i\ j)(i::q(j::v'))(j::z))v.
\]
The inner recursive value z is the required same-first-word/shorter-
second-word call. Repeated variable indices are retained, not deduplicated.
\end{hbdefinition}

\begin{hblemma}{thm:p12-har-poly-word-order}{proved=annex:p12.har.poly-word-order}{All word-code cases and the strict total ordering}
Weq and wlt satisfy Bit; weq u v=one iff u=v in primitive equality
on W, and wlt u v=one iff u<lex v. Their O codes are equivalent
to the respective negations. Lex order is irreflexive, transitive,
asymmetric and trichotomous. No word-equality oracle is assumed.
\end{hblemma}

\begin{hbdefinition}{def:p12-har-poly-invariants}{The concrete sparse-data invariants}
Use the inductive element-occurrence relation on W (head and tail
constructors) and row-occurrence relation on P, again with head/tail
constructors. \(\mathrm{Sorted}_W\) has the empty and singleton cases, and
\(\mathrm{Sorted}_W\)(i::j::w) precisely when i<=j and \(\mathrm{Sorted}_W\)(j::w).
Good epsilon holds. Good(row c w \# p) precisely requires c not=O,
\(\mathrm{Sorted}_W\) w, Good p and, for every row c' v occurring in p, v<lex w.
These are Prop predicates with the displayed generated proof branches;
they are not proof fields consumed by the program.
\end{hbdefinition}

\begin{hbdefinition}{def:p12-har-poly-semantics}{Actual coefficient, word and polynomial interpretations}
In the fixed commutative semiring and valuation put
\[
\begin{aligned}
 \chi&:=\mathsf{nat\_rec}\ 0_A
                (\lambda k:\mathsf{nat}.\lambda a:A.a+_A1_A),\\
 V_\nu&:=\mathrm{Rec}_W1_A
       (\lambda i:\mathsf{nat}.\lambda w:W.\lambda a:A.\nu i\times_A a),\\
 B_\nu(\mathsf{row}\ c\ w)&:=\chi c\times_A V_\nu w,\\
 L_\nu&:=\mathrm{Rec}_P0_A
       (\lambda r:R.\lambda p:P.\lambda a:A.B_\nu r+_A a).
\end{aligned}
\]
The row interpretation is its \(\mathrm{Rec}_R\) branch; targets are the fixed A.
These are actual functions, not interpretations postulated to satisfy
the normalizer's desired conclusion.
\end{hbdefinition}

\begin{hblemma}{thm:p12-har-poly-coefficients}{proved=annex:p12.har.poly-coefficients}{Every coefficient homomorphism clause}
\(\chi(0)=0_A\), \(\chi(\mathsf{one})=1_A\),
\(\chi(c+d)=\chi c+_A\chi d\), and
\(\chi(cd)=\chi c\times_A\chi d\).
In the Nat instance \(\chi(n)=n\). The proofs use the declared semiring
laws and generated nat induction, not an assumed homomorphism.
\end{hblemma}

\begin{hblemma}{thm:p12-har-poly-word-merge}{proved=annex:p12.har.poly-word-merge}{The actual merge preserves sortedness and every variable occurrence}
MW preserves \(\mathrm{Sorted}_W\) on two sorted inputs. An index occurs in MW u v
iff it occurs in u or v. In any declared commutative-semiring interpretation,
\(V_\nu(\mathsf{MW}(u,v))=V_\nu(u)\times_A V_\nu(v)\),
where V is the word interpretation defined above. The semantic equation
holds even on unsorted words; commutativity is an explicit premise.
\end{hblemma}

\begin{hbdefinition}{def:p12-har-poly-insert}{Actual insertion, sum, scale and product}
Put mk(c,w,p):=\(I_P\)(eqb c O)p(row c w \# p). Define
\[
 \mathsf{Ins}(c,w,p):=
 I_P(\mathsf{eqb}\ c\ O)p
  \bigl(\mathrm{Rec}_P(\mathsf{mk}(c,w,\epsilon))f_{c,w}\ p\bigr),
\]
where the branch on head row d v, tail s and recursive value z:P is
\[
 f_{c,w}(\mathsf{row}\ d\ v,s,z):=
 I_P(\mathsf{weq}(w,v))\mathsf{mk}(c+d,w,s)
  \bigl(I_P(\mathsf{wlt}(v,w))
    (\mathsf{row}\ c\ w\mathbin{\#}(\mathsf{row}\ d\ v\mathbin{\#}s))
    (\mathsf{row}\ d\ v\mathbin{\#}z)\bigr).
\]
The branch's row pattern is the single \(\mathrm{Rec}_R\)/case branch, not a
destructive mutation. Zero input coefficients return p; an equal key
combines coefficients, a larger new key is prepended, and a smaller
one recurses through the tail. The mk zero test is present also at
the combination branch.
\[
\begin{aligned}
 \mathsf{Sum}(p,q)&:=\mathrm{Rec}_P q
  (\lambda r:R.\lambda s:P.\lambda z:P.
        \mathrm{case}_R(\lambda c:\mathsf{nat}.\lambda w:W.
                            \mathsf{Ins}(c,w,z))r)\ p,\\
 \mathsf{Scale}(c,w,p)&:=\mathrm{Rec}_P\epsilon
  (\lambda r:R.\lambda s:P.\lambda z:P.
   \mathrm{case}_R(\lambda d:\mathsf{nat}.\lambda v:W.
              \mathsf{Ins}(cd,\mathsf{MW}(w,v),z))r)\ p,\\
 \mathsf{Prod}(p,q)&:=\mathrm{Rec}_P\epsilon
  (\lambda r:R.\lambda s:P.\lambda z:P.
   \mathrm{case}_R(\lambda c:\mathsf{nat}.\lambda w:W.
              \mathsf{Sum}(\mathsf{Scale}(c,w,q),z))r)\ p.
\end{aligned}
\]
Every recursive target here is P; all source syntax is expanded into
the displayed ordinary data-recursion applications.
\end{hbdefinition}





\begin{hblemma}{thm:p12-har-poly-insertion}{proved=annex:p12.har.poly-insertion}{Every insertion branch, key order and value}
For Good p and \(\mathrm{Sorted}_W\) w, Ins(c,w,p) is Good, including c=O.
Every output key is either w or an input key. For arbitrary raw p,w,c,
\[
 L_\nu(\mathsf{Ins}(c,w,p))=
           (\chi c\times_A V_\nu w)+_A L_\nu p.
\]
Inserting a positive row whose key is larger than all input keys gives
that row followed by the unchanged input, in primitive book equality.
\end{hblemma}

\begin{hblemma}{thm:p12-har-poly-operations}{proved=annex:p12.har.poly-operations}{The complete computed sum/product and all zero cases}
Sum and Prod preserve Good on Good inputs; Scale(c,w,p) preserves
Good when w is sorted and p is Good. Their actual interpreted values are
\[
\begin{aligned}
 L_\nu(\mathsf{Sum}(p,q))&=L_\nu p+_A L_\nu q,\\
 L_\nu(\mathsf{Scale}(c,w,p))&=(\chi c\times_A V_\nu w)\times_A L_\nu p,\\
 L_\nu(\mathsf{Prod}(p,q))&=L_\nu p\times_A L_\nu q.
\end{aligned}
\]
The semantic equations hold on raw data. Sum(epsilon,q)=q;
Prod(epsilon,q) and Prod(p,epsilon) return epsilon; Scale(0,w,p)
returns epsilon. Zero coefficients are dropped rather than retained
as false nonzero rows.
\end{hblemma}

\begin{hbdefinition}{def:p12-har-poly-normalizer}{All five actual normalization branches}
Write single(c,w):=mk(c,w,epsilon). The actual data normalization is
\[
 N:=\mathsf{fold}_T^P\bigl(
 \epsilon,\mathsf{single}(\mathsf{one},[]),
 \lambda i:\mathsf{nat}.\mathsf{single}(\mathsf{one},i::[]),
 \lambda s,t:T.\lambda p,q:P.\mathsf{Sum}(p,q),
 \lambda s,t:T.\lambda p,q:P.\mathsf{Prod}(p,q)\bigr).
\]
Thus N:T-to-P is an actual fold, not the assertion that some normal
polynomial exists.
\end{hbdefinition}

\begin{hbtheorem}{thm:p12-har-poly-normalization}{proved=annex:p12.har.poly-normalization}{The complete five-branch normal-form and universal value proof}
For every t:T, Good(Nt), and for every valid open valuation in the
declared commutative semiring,
\(L_\nu(Nt)=E_\nu^A t\).
The program is valuation-independent; the preservation theorem is not.
\end{hbtheorem}



\begin{hbdefinition}{def:p12-har-poly-quotation}{Actual quotation-back and G}
The four concrete quotation recursions are
\[
\begin{aligned}
 \mathsf{lit}&:=\mathsf{nat\_rec}\ Z
    (\lambda k:\mathsf{nat}.\lambda t:T.
        I_T(\mathsf{eqb}\ k\ O)U(\mathsf{Add}\ t\ U)),\\
 Q_W&:=\mathrm{Rec}_W U
    (\lambda i:\mathsf{nat}.\lambda w:W.\lambda t:T.
       I_T(\mathsf{empty}_W w)(\mathsf{Var}\ i)
                    (\mathsf{Mul}(\mathsf{Var}\ i)t)),\\
 Q_R(\mathsf{row}\ c\ w)&:=
   I_T(\mathsf{empty}_W w)(\mathsf{lit}\ c)
      (I_T(\mathsf{eqb}\ c\ \mathsf{one})(Q_Ww)
                          (\mathsf{Mul}(\mathsf{lit}\ c)(Q_Ww))),\\
 Q_P&:=\mathrm{Rec}_P Z
   (\lambda r:R.\lambda p:P.\lambda t:T.
        I_T(\mathsf{empty}_P p)(Q_Rr)(\mathsf{Add}(Q_Rr)t)),\\
 G&:=\lambda t:T.Q_P(Nt).
\end{aligned}
\]
\(Q_R\) is the displayed single \(\mathrm{Rec}_R\) branch. All four targets are T.
The source's \(x^2\) and \(2x\) are printing sugar for Mul X X and
Mul(Add U U)X; no additional constructor is inserted into T.
\end{hbdefinition}

\begin{hblemma}{thm:p12-har-poly-quotation}{proved=annex:p12.har.poly-quotation}{All quotation branches and the actual value bridge}
For every raw datum,
\(E_\nu^A(\mathsf{lit}\ c)=\chi c\),
\(E_\nu^A(Q_W w)=V_\nu w\),
\(E_\nu^A(Q_R r)=B_\nu r\), and
\(E_\nu^A(Q_P p)=L_\nu p\).
Each zero, empty, singleton, coefficient-one and other-coefficient case
is constructed in the annex.
\end{hblemma}

\begin{hbtheorem}{thm:p12-har-poly-preservation}{proved=annex:p12.har.poly-preservation}{G's full universal open-valuation preservation}
For every fixed declared commutative semiring, every valid-context
\(\nu:\mathsf{nat}\to A\), and every t:T,
\[
 E_\nu^A(Gt)=E_\nu^A t.
\]
In particular this supplies the source Nat obligation at every open
natural-number valuation. No equality of valuation functions or closure
of the input expression is assumed.
\end{hbtheorem}

\begin{hbtheorem}{thm:p12-har-poly-retraction}{proved=annex:p12.har.poly-retraction}{The actual normal-form retraction and idempotence}
For Good p, \(N(Q_Pp)=p\) in primitive equality on P.
Consequently G(Gt)=Gt in primitive equality on T, and every output is
the declared quotation of Good sparse data. These are book equalities,
not a claim of definitional idempotence on opaque open indices or of
semantic completeness for arbitrary semiring models.
\end{hbtheorem}

\begin{hbtheorem}{thm:p12-har-poly-example}{proved=annex:p12.har.poly-example}{The source square-expansion example and its precise output convention}
Put X:=Var O. The actual computation is
\[
\begin{aligned}
 G(\mathsf{Mul}(\mathsf{Add}\ X\ U)(\mathsf{Add}\ X\ U))
 \twoheadrightarrow_{\beta\iota\delta}
 \mathsf{Add}(\mathsf{Mul}\ X\ X)
  (\mathsf{Add}(\mathsf{Mul}(\mathsf{Add}\ U\ U)X)U).
\end{aligned}
\]
This is \(x^2\)+(2x+1) under the declared right-associated quotation.
The left-associated source reading (\(x^2\)+2x)+1 has the same interpretation
by the explicit addition-associativity law. Literal AST identity with
that alternate parsing is not asserted. This is a nontrivial symbolic
output of the constructed G, not identity G or only a numerical example
\cite[pp.~1221--1222]{BarendregtGeuvers2001}.
\end{hbtheorem}

\begin{proofboundary}{def:p12-har-poly-residual}{What the actual normalizer does not identify}
The program is a sparse formal natural-coefficient polynomial normalizer
for the declared five-constructor grammar. Preservation for a non-Nat
interpretation requires its actual commutative-semiring laws; canonical
data do not imply semantic completeness for a model with extra equations.
Quotation's association convention is explicit rather than an assertion
of a parser-dependent source AST identity. The distinct generic computation
construction is given in
\cref{thm:p12-har-prcompiler-representation,thm:p12-har-prcompiler-ok,thm:p12-har-prcompiler-predicate};
polynomial normalization alone does not supply it or arbitrary proposition
reification. Pocklington/Fermat and large-number certificates, differentiation and
continuity extensions, arithmetic coding and global extension metatheory
remain distinct source obligations. No kernel replay is claimed here.
\end{proofboundary}
